\title{xLSTM: Extended Long Short-Term Memory}

\begin{abstract}
During the 1990s, the Long Short-Term Memory (LSTM) architecture was defined by the introduction of the constant error carousel and gating mechanisms. LSTMs have since demonstrated enduring utility and have underpinned numerous advancements in deep learning—most notably, they were foundational for early Large Language Models (LLMs). The emergence of the Transformer architecture, characterized by its scalable and highly parallelizable self-attention mechanism, subsequently shifted the paradigm by significantly outperforming LSTMs at scale. This work re-examines language modeling by scaling LSTM-based systems up to billions of parameters, leveraging innovations derived from the latest LLM methodologies, while also directly addressing and alleviating well-known weaknesses of LSTMs. Our contributions are twofold: we develop an exponential gating mechanism accompanied by tailored normalization and stabilization protocols; additionally, we re-engineer the core LSTM memory, resulting in: (i) the sLSTM, which utilizes scalar memory and updates together with an enhanced memory mixing strategy, and (ii) the mLSTM, an entirely parallelizable variant based on matrix memory and a covariance-based update schema. Embedding these improved LSTM forms within the structure of residual blocks, we propose xLSTM blocks, which are composed into deep xLSTM model architectures via residual stacking. The combination of exponential gating and restructured memory endows xLSTM models with strong empirical performance and scalability, often surpassing leading Transformer and State Space Model approaches.\\
Code available at: \url{https://github.com/NX-AI/xlstm}
\end{abstract}

\section{Introduction}
\label{sec:intro}

The Long Short-Term Memory (LSTM) ideas~\citep{Hochreiter:91,Hochreiter:97nips,Hochreiter:97}, i.e., the constant error carousel and gating, were introduced to overcome the vanishing gradient problem of recurrent neural networks~\citep{Hochreiter:91,Hochreiter:00book}:

\begin{align}
\label{eq:lstm_idea}
  c_t \ = \ f_t \ c_{t-1} \ + \ 
          i_t \  z_t \ , \quad  
          h_t \ = \ o_t \ \psi({c_t}) \ . 
\end{align}

The constant error carousel operates by incrementally adjusting the cell state $c_{t-1}$ (green) with the addition of cell inputs $z_t$, which are regulated by sigmoid gating mechanisms (blue). The input gate $i_t$ and forget gate $f_t$ are responsible for determining the extent of this adjustment, whereas the output gate $o_t$ dictates what portion of the cell’s memory is revealed in the hidden state $h_t$. Prior to output, the cell state is processed via the squashing or normalization function $\psi$, after which the output gate controls the final hidden state.

LSTMs have achieved remarkable success across a range of application areas \citep{Hochreiter:01,Hochreiter:07,Schmidhuber:15}, dominating text generation tasks prior to the rise of Transformers in 2017~\citep{Vaswani:17}. Numerous studies have confirmed the capabilities of LSTMs on sequential tasks such as text generation~\citep{Graves:13,Karpathy:15blog}, handwriting synthesis~\citep{Graves:13}, sequence-to-sequence machine translation~\citep{Sutskever:14nips}, program evaluation~\citep{Zaremba:14arxiv}, automatic image captioning~\citep{Karpathy:15,Hossain:19}, source code generation~\citep{Karpathy:15blog}, hydrological tasks including rainfall-runoff modeling~\citep{Kratzert:18,Kratzert:19}, and the development of hydrological warning systems~\citep{Nearing:24}. Within reinforcement learning, LSTMs have established themselves as the leading architecture for modeling sequences, as evidenced by their incorporation in models such as AlphaStar for StarCraft II~\citep{Vinyals:17short}, OpenAI Five for Dota 2~\citep{OpenAI:19}, and the magnetic controller in nuclear fusion research~\citep{Degrave:22short}. Their strength lies in their proficiency at capturing abstract representations and storing salient semantic information within their memory cells~\citep{Karpathy:15blog}, a faculty illustrated by discoveries of number and syntax neurons~\citep{Lakretz:19}, linguistic features~\citep{Bau:19}, and sentiment neurons~\citep{Radford:17arxiv}. LSTMs continue to be foundational in significant contemporary uses~\citep{Degrave:22short,Nearing:24}, consistently demonstrating their enduring relevance.

Although LSTMs have achieved significant progress, they are hindered by three principal shortcomings: (i) Inability to reconsider prior storage choices. This shortcoming is illustrated by the 
{\it Nearest Neighbor Search} task (refer also to Appendix~\ref{sec:appNearestNeighborSearch}). Given a reference vector, the sequence must be traversed in order to identify the most similar vector, whose associated value is produced at the sequence’s end. The left-hand panel of Figure~\ref{fig:lstmProblems} presents the average squared error on this task. LSTMs exhibit difficulty in replacing a stored value when a better match occurs later, whereas our proposed xLSTM addresses this through exponential gating. (ii) Constrained memory capacity requiring all information to be compressed into a single scalar state per cell. This is exemplified using the {\it Rare Token Prediction} benchmark. On the right in Figure~\ref{fig:lstmProblems}, token prediction perplexity on the Wikitext-103 dataset~\citep{Merity:17} is displayed, partitioned by token occurrence rate. LSTM exhibits suboptimal performance on infrequent tokens due to insufficient storage capacity. The xLSTM overcomes this by introducing a matrix-based memory. (iii) Inability to parallelize across time due to memory entanglement, where transitions between hidden states enforce sequential computation.

These drawbacks of LSTM architectures facilitated the development of Transformers~\citep{Vaswani:17} within the field of language modeling. An important question arises: What modeling outcomes are attainable if these LSTM deficiencies are resolved and LSTMs are upscaled to rival contemporary Large Language Model sizes?

\section{Extended Long Short-Term Memory}
\label{sec:xLSTM}

In response to the constraints of standard LSTMs, Extended Long Short-Term Memory (xLSTM) introduces two substantial modifications to the core LSTM framework of Equation~\eqref{eq:lstm_idea}. These enhancements—namely exponential gating and innovative memory structures—expand the LSTM repertoire with two novel variants: (i) sLSTM, presented in Section~\ref{sec:sLSTM}, which utilizes scalar memory, scalar updates, and integrates memory mixing, and (ii) mLSTM, detailed in Section~\ref{sec:mLSTM}, which features a matrix-based memory and implements a covariance (outer product) updating mechanism suitable for full parallel execution. The exponential gating scheme is fundamental to both sLSTM and mLSTM, promoting improved performance. Unlike sLSTM, the mLSTM excludes the use of memory mixing—or hidden-hidden recurrent dynamics—to support parallel computation. Both structures, sLSTM and mLSTM, are adaptable to architectures employing multiple memory cells; in sLSTM, this comes with the added complexity of memory mixing between cells. Additionally, sLSTM may include multiple heads with intra-head cell-level memory mixing, while avoiding mixing between different heads—constituting a novel form of memory interaction introduced through exponential gating. For mLSTM, multiple heads and multiple cells serve an equivalent role.

The integration of these enhanced LSTM models into residual block frameworks produces xLSTM blocks, as explained in Section~\ref{sec:xLSTMblock}. By stacking these blocks with residual connections, xLSTM architectures can be constructed, as outlined in Section~\ref{sec:xLSTMarchitecure}. Figure~\ref{fig:xlstm_sketch} provides an overview of the xLSTM architecture and its constituent elements.

\subsection{Review of the Long Short-Term Memory}
\label{sec:orgLSTM}

The foundational LSTM concept~\citep{Hochreiter:91,Hochreiter:97nips,Hochreiter:97} proposed a scalar memory cell serving as the core mechanism for processing and storing information. This architecture mitigates the issue of vanishing gradients~\citep{Hochreiter:91,Hochreiter:00book} by implementing the constant error carousel, specifically through the updating of the cell state. Within the memory cell, three distinct gates are incorporated: input, output, and forget gates. The forget gate was subsequently added to the architecture by \citet{Gers:00}. The LSTM memory cell updates at each time step $t$ are governed by the following set of rules:

\begin{align}
c_t \ &= \  f_t \ c_{t-1} \ + \ i_t \ z_t &  & &\text{cell state} \\
h_t  \ &= \ o_t \ \tilde{h}_t \ , & \tilde{h}_t \ &= \ \psi \left( c_t\right)
  &\text{hidden state} \\
z_t \ &= \ \varphi \left( \tilde{z}_t \right) \ , 
  &\tilde{z}_t \ &=  \ \mathbf{w}^\top_{z} \ \mathbf{x}_t \ + \
  r_{z}  h_{t-1} \ + \  b_{z} \ \
  &\text{cell input} \\
i_t \ &= \ \sigma \left( \tilde{i}_t  \right) \ , 
  &\tilde{i}_t \ &= \ \mathbf{w}^\top_{i} \ \mathbf{x}_t \ + \
  r_{i}  \ h_{t-1} \ + \  b_{i} \ \
  &\text{input gate} \\
f_t \ &= \ \sigma \left(  \tilde{f}_t \right) \ , 
  &\tilde{f}_t \ &= \ \mathbf{w}^\top_{f} \ \mathbf{x}_t  \ + \
  r_{f}  \ h_{t-1} \ + \  b_{f} \ \
  &\text{forget gate} \\
o_t \ &= \ \sigma \left( \tilde{o}_t \right) \ , 
  &\tilde{o}_t  \ &= \ \mathbf{w}^\top_{o} \ \mathbf{x}_t \ + \
  r_{o}  \ h_{t-1} \ + \  b_{o} \ \
  &\text{output gate} 
\end{align}

The vectors $\mathbf{w}_{z}$, $\mathbf{w}_{i}$, $\mathbf{w}_{f}$, and $\mathbf{w}_{o}$ serve as the input weight vectors connecting the input $\mathbf{x}_t$ to the cell input, input gate, forget gate, and output gate, respectively. The recurrent weights, denoted as $r_{z}$, $r_{i}$, $r_{f}$, and $r_{o}$, represent the connections between the previous hidden state $h_{t-1}$ and each of the corresponding gates and the cell input. Bias terms are denoted by $b_{z}$, $b_{i}$, $b_{f}$, and $b_{o}$, each associated with their respective components. The functions $\varphi$ and $\psi$ indicate the activation functions applied to the cell input and hidden state, respectively (most often $\tanh$ is utilized). The function $\psi$ specifically operates to bound the values of the cell state, preventing it from growing without limit. Each gate applies a sigmoid function as its activation: $\sigma \left( x \right) = 1/(1+\exp(-x))$. Later developments grouped several scalar memory cells $\mathbf{c}_t \in \mathbb{R}^{d}$ into a vector form, enabling the introduction of recurrent weight matrices $\mathbf{R} \in \mathbb{R}^{d \times d}$ that facilitate interactions among the outputs of different memory cells~\citep{Greff:15}; refer to Appendix~\ref{sec:apporgLSTM} for additional information. Ablation experiments have demonstrated that all elements of the memory cell architecture are essential for its performance~\citep{Greff:15}.

\subsection{sLSTM}
\label{sec:sLSTM}

In order to enable LSTMs to reconsider their memory retention choices, we incorporate exponential gates (depicted in red), normalization, and stabilization mechanisms. Specifically, the input and forget gates are augmented to use exponential activation functions. For normalization purposes, we define a normalizer state that accumulates the sum of the input gate values multiplied by all subsequent forget gate values.

The scalar sLSTM forward pass is:

\begin{align}
\label{eq:slstmforward}
c_t \ &= \  f_t \ c_{t-1} \ + \ i_t \ z_t & & 
 &\text{cell state} \\
n_t \ &= \  f_t \ n_{t-1} \ + \ i_t  & & 
 &\text{normalizer state} \\
h_t \ &= \ o_t \ \tilde{h}_t \ ,
  &\tilde{h}_t \ &= \ c_t / n_t
&\text{hidden state} \\
z_t \ &= \ \varphi \left( \tilde{z}_t \right) \ , 
  &\tilde{z}_t \ &=  \ \mathbf{w}^\top_{z} \ \mathbf{x}_t \ + \
  r_{z}  h_{t-1} \ + \  b_{z}
  &\text{cell input} \\
i_t \ &= \ \!\exp \left( \tilde{i}_t  \right) \ , 
  &\tilde{i}_t \ &= \ \mathbf{w}^\top_{i} \ \mathbf{x}_t \ + \
  r_{i}  \ h_{t-1} \ + \  b_{i}
  &\text{input gate} \\
f_t \ &= \ \sigma \left(  \tilde{f}_t \right) \ \text{OR} \
\!\exp \left(  \tilde{f}_t \right) \ ,
  &\tilde{f}_t \ &= \ \mathbf{w}^\top_{f} \ \mathbf{x}_t  \ + \
  r_{f}  \ h_{t-1} \ + \  b_{f}
  &\text{forget gate} \\
o_t \ &= \ \sigma \left( \tilde{o}_t \right) \ , 
  &\tilde{o}_t  \ &= \ \mathbf{w}^\top_{o} \ \mathbf{x}_t \ + \
  r_{o}  \ h_{t-1} \ + \  b_{o}
  &\text{output gate} 
\end{align}

We transfer the original LSTM gating techniques, i.e., input- and/or hidden-dependent gating plus bias term, to the new architectures. Exponential activation functions can lead to large values that cause overflows. Therefore, we stabilize gates with an additional state \mbox{$m_t$~\citep{Milakov:18arxiv}}:

\begin{align}
\label{eq:slstmstabil}
m_t \ &= \ \max \left( \log ( f_t ) + m_{t-1} , \log ( i_t ) \right) &\text{stabilizer state} \\
i'_t \ &= \ \!\exp \left( \log \left ( i_t \right) - m_t \right) \ = \ \!\exp \left( \tilde{i}_t - m_t \right) \ \, 
  &\text{stabil. input gate} \\
f'_t \ &= \ \!\exp \left( \log \left( f_t \right) + m_{t-1} - m_t \right) \ \, 
  &\text{stabil. forget gate}
\end{align}

As demonstrated in Appendix~\ref{sec:appsLSTM}, substituting $f_t$ with $f'_t$ and $i_t$ with $i'_t$ during the forward computation does not alter either the network’s overall output or the gradients of the loss relative to its parameters.

\paragraph{New Memory Mixing.} Like the conventional LSTM, sLSTM is capable of incorporating several memory cells (refer to Appendix~\ref{sec:appsLSTM}). This multi-cell arrangement allows the mixing of memory states through recurrent connections denoted as $\mathbf{R}_{\mathbf{z}}$, $\mathbf{R}_{\mathbf{i}}$, $\mathbf{R}_{\mathbf{f}}$, $\mathbf{R}_{\mathbf{o}}$, which transmit information from the hidden state vector $\mathbf{h}$ to the memory cell input $\mathbf{z}$ and to the respective gates $\mathbf{i}$, $\mathbf{f}$, and $\mathbf{o}$. The exponential gating mechanism introduces a new dynamic in memory integration. The revised sLSTM architecture can support several heads, wherein memory mixing is permitted within each head, but not between heads. The combination of multiple heads and exponential gating thus provides a novel approach to memory mixing in sLSTM.

\subsection{mLSTM}
\label{sec:mLSTM}

To improve the memory capacity of LSTMs, we expand the memory cell representation from a single scalar $c \in \mathbb{R}$ to a matrix $\mathbf{C} \in \mathbb{R}^{d \times d}$. Consequently, the retrieval operation employs matrix multiplication. At each time step $t$, the goal is to store a pair of vectors—a key $\mathbf{k}_t \in \mathbb{R}^d$ and a value $\mathbf{v}_t \in \mathbb{R}^d$—adopting the nomenclature from Transformer models. At some later time $t + \tau$, a query vector $\mathbf{q}_{t+\tau} \in \mathbb{R}^d$ is used to access the previously stored value $\mathbf{v}_t$. This framework mirrors the setup found in Bidirectional Associative Memories (BAMs) \citep{Kohonen:72,Anderson:72,Nakano:72,Anderson:77}.

The covariance update rule~\citep{Sejnowski:77,Dayan:91} for storing a key-value pair is

\begin{align}
\mathbf{C}_t \ &= \ \mathbf{C}_{t-1} \ + \ \mathbf{v}_{t} \ \mathbf{k}_t^\top \ .
\end{align}

We apply layer normalization to the inputs prior to their projection into key and value spaces, ensuring these vectors possess zero mean. According to~\citep{Dayan:91}, the covariance update rule provides an optimal means for maximizing the separability among the retrieved binary representations, which directly corresponds to achieving a maximal signal-to-noise ratio. This separation can be increased if one restricts retrieval operations to pairwise relations and is willing to accept the quadratic computational cost characteristic of attention mechanisms~\citep{Krotov:16,Krotov:17,Ramsauer:21}. The aforementioned covariance update strategy aligns with the Fast Weight Programmers~\citep{Schmidhuber:92ncfastweights,Schlag:21}, which were subsequently adapted to include a constant decay rate applied to $\mathbf{C}_{t-1}$ alongside a constant learning rate applied to 
$\mathbf{v}_{t} \mathbf{k}_t ^\top$~\citep{Ba:16}. Guided by this perspective, we embed the covariance update rule into the architecture of LSTMs, interpreting the forget gate as analogous to the decay rate, the input gate as the learning rate, and viewing the output gate as responsible for modulating the magnitude of the retrieved output vector.

For this matrix memory, the normalizer state is the weighted sum of key vectors, where each key vector is weighted by the input gate and all future forget gates. Again, the normalizer state keeps record of the strength of the gates. Since the dot product between query and normalizer state can be close to zero, we use the absolute value of this dot product and lower bound it by a threshold (typically 1.0) as done previously~\citep{Sun:23arxiv}. The mLSTM forward pass is:

\begin{align}
\label{eq:mlstm_recurrent_begin}
\mathbf{C}_t \ &= \  f_t \ \mathbf{C}_{t-1} \ + \ 
  i_t \ \mathbf{v}_t \ \mathbf{k}_t^\top & &  &\text{cell state} \\
\mathbf{n}_t \ &= \  f_t \ \mathbf{n}_{t-1} \ + \ 
  i_t \ \mathbf{k}_t & &  &\text{normalizer state} \\
\mathbf{h}_t  \ &= \ \mathbf{o}_t \ \odot \ \tilde{\mathbf{h}}_t
\ , \qquad \qquad 
\tilde{\mathbf{h}}_t \ = \   \mathbf{C}_t \mathbf{q}_t \ / \ 
  \max \left\{ |\mathbf{n}_t^\top \mathbf{q}_t|, 1 \right\} 
   &&&\text{hidden state} \\
\mathbf{q}_t \ &= \ \mathbf{W}_q \ \mathbf{x}_t \ + \ \mathbf{b}_q  & &  &\text{query input} \\
\mathbf{k}_t \ &= \ \frac{1}{\sqrt{d}} \mathbf{W}_k \ \mathbf{x}_t \ + \ \mathbf{b}_k  & &  &\text{key input} \\
\mathbf{v}_t \ &= \ \mathbf{W}_v \ \mathbf{x}_t \ + \ \mathbf{b}_v  & &  &\text{value input} \\
i_t \ &= \ \!\exp \!  \! \left( \tilde{i}_t  \right) \ , 
 \qquad \qquad \ \ \ \ \,
  \tilde{i}_t \ = \ \mathbf{w}^\top_{i} \ \mathbf{x}_t \ + \  b_{i} \ \ \ 
  &&&\text{input gate} \\
f_t \ &= \ \sigma \!  \left(  \tilde{f}_t \right) \ \text{OR} \ \!\exp \!  \! \left(  \tilde{f}_t \right) , \ \ \: \!
\tilde{f}_t \ = \ \mathbf{w}^\top_{f} \ \mathbf{x}_t  \ + \
  b_{f}
  &&& \text{forget gate} \\
\label{eq:mlstm_recurrent_end}
\mathbf{o}_t \ &= \ \sigma \left( \tilde{\mathbf{o}}_t \right) \ , \qquad \qquad \quad \ \ \ \,
  \tilde{\mathbf{o}}_t  \ = \ \mathbf{W}_{\mathbf{o}} \ \mathbf{x}_t \ + \
  \mathbf{b}_{\mathbf{o}}
  &&&\text{output gate} 
\end{align}

The mLSTM architecture is capable of supporting several memory cells, similar to the traditional LSTM. In the case of mLSTM, having multiple cells functions equivalently to incorporating multiple heads, given the absence of any interaction between memory components. To address potential instability arising from the exponential gating mechanisms in mLSTM, we apply the same stabilization strategies utilized for sLSTM (refer to Equation~\ref{eq:slstmstabil}). Because mLSTM does not involve inter-cell memory mixing, the associated recurrence relation permits a reformulation amenable to parallel computation. Additional information is provided in Appendix~\ref{sec:appmLSTM}.

\subsection{xLSTM Architecture}
\label{sec:xLSTMarch}

\paragraph{xLSTM Blocks.}
\label{sec:xLSTMblock}
An xLSTM block is designed to encode past information in a non-linear fashion within a high-dimensional space, thereby enabling finer separation of different histories or contextual sequences. The capacity to distinguish between distinct histories is essential for making accurate predictions about subsequent elements in a sequence, such as the following token. This approach draws upon Cover’s Theorem~\citep{Cover:65}, which posits that patterns embedded through non-linear transformations into higher-dimensional spaces become more amenable to linear separation relative to their lower-dimensional counterparts. We analyze two types of residual block structures: (i) A residual block utilizing post up-projection (similar to that in Transformers), where the mechanism first performs a non-linear encoding of the past within the original feature space, then employs a linear mapping to increase dimensionality, follows this with a non-linear activation, and finally projects back linearly to the initial space; this is illustrated in the left panel of Figure~\ref{fig:Backbones} and detailed further in the third column of Figure~\ref{fig:xlstm_sketch}. For an expanded depiction, see Figure~\ref{fig:appsLSTM_detailed} in the appendix. (ii) Alternatively, a residual block incorporating pre up-projection (resembling State Space Models), which first projects the data linearly into a higher-dimensional space,

It performs a nonlinear abstraction of previous information within a high-dimensional feature space, which is subsequently projected back to the original dimensionality through a linear transformation. When an xLSTM block employs an sLSTM, the architecture incorporates a post up-projection module. In contrast, for an xLSTM block based on an mLSTM, we opt for a pre up-projection module due to the expanded memory storage achieved in higher dimensions. Further clarification is provided in the leftmost panel of Figure~\ref{fig:Backbones}, the third column in Figure~\ref{fig:xlstm_sketch}, as well as Figure~\ref{fig:appsLSTM_detailed} in the appendix.

\paragraph{xLSTM Architecture. }
\label{sec:xLSTMarchitecure}
The xLSTM model architecture is formed by residual stacking of foundational modules~\citep{Srivastava:15,He:16}. Consistent with prevailing practices in Large Language Models, we utilize a pre-LayerNorm~\citep{Ba:16b} residual structure as the backbone. The rightmost panel in Figure~\ref{fig:xlstm_sketch} provides an illustration of this arrangement.

\subsection{Memory and Speed Considerations}

In contrast to Transformers, xLSTM networks feature linear computational scaling and maintain constant memory requirements as sequence length increases. The compressive memory mechanism of xLSTM makes it highly appropriate for industrial contexts and edge-device deployment.

Although mLSTM’s memory operates without needing trainable parameters, its $d \times d$ matrix memory and $d \times d$ update scheme render it computationally intensive. This represents a compromise, increasing computational demands to gain higher memory capacity. However, the associated operations are efficiently parallelizable on modern GPUs, and thus exert only a marginal influence on overall processing time.

Unlike mLSTM, whose architecture allows for parallelization akin to approaches such as FlashAttention~\citep{Dao:22,Dao:23} and GLA~\citep{Yang:23arxiv}, sLSTM cannot be parallelized due to the nature of its memory mixing through hidden-to-hidden interactions. To address this, we engineered a high-performance CUDA implementation with GPU register-level memory enhancements, resulting in sLSTM processing speeds that are typically only up to twice as slow as those observed for mLSTM.

\section{Related Work}
\label{sec:related}

\paragraph{Linear Attention.} A range of techniques have been proposed to mitigate the quadratic scaling of the Transformer's attention mechanism with respect to context length, thereby enabling attention to scale linearly. The Synthesizer model, for instance, constructs artificial attention scores independently of token-to-token relationships~\citep{Tay:20arxiv}. Linformer utilizes a low-rank matrix to represent self-attention and achieves a linear approximation~\citep{Wang:20arxiv}. Linear Transformer reformulates the attention mechanism to ensure linear computational complexity~\citep{Katharopoulos:20}. Performer leverages positive orthogonal random features to approximate the attention softmax in a linear manner~\citep{Choromanski:21}. Moreover, rapid convolutional operations have replaced traditional attention in methods such as Structured Global Convolution (SGConv)~\citep{Li:22} and Hyena Hierarchy~\citep{Poli:23}.

\paragraph{State Space Models.} State Space Models (SSMs) have gained considerable attention in recent times, attributed to their linear scaling with context length and impressive empirical results that rival those of Transformers. Initial advancements in this area include the introduction of the Structured State Space sequence model (S4)~\citep{Gu:21}. Subsequently, various extensions and modifications have been proposed, such as the Diagonal State Space (DSS) model~\citep{Gupta:22}, Gated State Space (GSS) models~\citep{Mehta:22}, S5 model~\citep{Smith:22}, Bidirectional Gated SSM (BiGS)~\citep{Wang:22}, H3 model~\citep{Fu:23}, and Mamba~\citep{Gu:24arxiv}.

\paragraph{Recurrent Neural Networks.} Recurrent Neural Networks (RNNs) have been proposed as an alternative to Transformer and attention mechanisms, mainly due to their linear scaling with input sequence length. Recent advances utilizing Deep Linear Recurrent Units (LRUs) have yielded strong outcomes in language modeling tasks~\citep{Orvieto:23, De:24arxiv}. Similarly, both Hierarchically Gated Linear RNN (HGRN)~\citep{Qin:23} and its successor HGRN2~\citep{Qin:24arxiv} have demonstrated substantial progress. Among RNN-based approaches, RWKV~\citep{Peng:23arxivshort,Peng:24arxivshort} stands out in large language modeling, achieving results comparable to Transformer models.

\paragraph{Gating.}
LSTM's introduction of gating stands as a foundational concept, subsequently revisited and reformulated in a variety of more recent models. This gating principle has been implemented in frameworks such as HGRN~\citep{Qin:23}, HGRN2~\citep{Qin:24arxiv}, Gated Linear Attention (GLA)~\citep{Yang:23arxiv}, Gated State Space (GSS) architectures~\citep{Mehta:22}, Bidirectional Gated SSM (BiGS)~\citep{Wang:22}, Moving Average Equipped Gated Attention (MEGA)~\citep{Ma:22}, as well as in architectures like RWKV~\citep{Peng:23arxivshort} and Mamba~\citep{Gu:24arxiv}.

\paragraph{Covariance Update Rule.}
To provide improved capacity for memory storage, the mLSTM cell has been augmented with a matrix memory utilizing a covariance-based update scheme. Comparable approaches that leverage such covariance updates include Fast Weight Programmers~\citep{Schmidhuber:92ncfastweights,Schlag:21}, RWKV-5 and RWKV-6~\citep{Peng:24arxivshort}, Retention~\citep{Sun:23arxiv}, Linear Transformer~\citep{Katharopoulos:20}, as well as HGRN2~\citep{Qin:24arxiv}.

\paragraph{Most Related.} The models most similar to xLSTM from a conceptual standpoint include Retention~\citep{Sun:23arxiv}, RWKV~\citep{Peng:23arxivshort,Peng:24arxivshort}, and HGRN2~\citep{Qin:24arxiv}. These frameworks incorporate matrix memory and/or employ gating mechanisms. Nevertheless, unlike the recently proposed sLSTM, these alternatives do not provide mechanisms for memory mixing. The capability for memory mixing is crucial for addressing state tracking tasks; as a result, LSTMs exhibit greater expressivity compared to State Space Models (SSMs) and Transformers~\citep{Merrill:24,Deletang:23}. State tracking functionality is essential for applications such as code execution or monitoring entities throughout extended narratives.

\paragraph{Residually Stacking Architectures.} The xLSTM architectures, similar to most state-of-the-art deep neural network models, are developed by stacking individual modules with residual connections~\citep{Srivastava:15,He:16}.

This approach has facilitated the development of deep convolutional networks~\citep{He:16} and has been foundational in the design of Transformers~\citep{Vaswani:17}. The Transformer paradigm is the driving innovation underlying the success of Large Language Models (LLMs) such as GPT-3~\citep{Brown:20short}, ChatGPT~\citep{Schulman:22short}, GPT-4~\citep{Achiam:23gpt4short}, Megatron-LM~\citep{Shoeybi:19}, Gopher~\citep{Rae:21short}, ERNIE 3.0 Titan~\citep{Wang:21arxivshort}, GLaM~\citep{Du:21glamshort}, Chinese M6~\citep{Lin:21}, multilingual AlexaTM 20B~\citep{Soltan:22}, OPT~\citep{Zhang:22opt}, Chinchilla~\citep{Hoffmann:22short}, BLOOM~\citep{Scao:22arxivshort}, GLM-130B~\citep{Zeng:22arxivshort}, LaMDA~\citep{Thoppilan:22short}, PaLM~\citep{Chowdhery:22short}, Llama~\citep{Touvron:23arxiv}, and Gemini~\citep{GeminiTeam:23arxiv,Reid:24arxiv}.

\section{Experiments}
\label{sec:Exp}

In this section, we carry out empirical assessments of xLSTM, emphasizing its application in language modeling and benchmarking it against established approaches. Section~\ref{sec:ExpSyntethic} explores xLSTM’s strengths using a range of synthetic tasks. 
Section~\ref{sec:ExpComparison} presents a comparative analysis of validation perplexity across several contemporary language modeling techniques, all trained with 15B tokens sourced from SlimPajama~\citep{Soboleva:23}. Here, we also include ablation studies specific to xLSTM. Furthermore, we examine the scaling properties of these methods, following methodologies outlined in \citet{Kaplan:20} and \citet{Brown:20short}.  
In Section~\ref{sec:ExpLanguage}, 
we pursue an extended language modeling investigation. In this experiment, both xLSTM and the top-performing techniques identified in Section~\ref{sec:ExpComparison} are trained on 300B tokens drawn from SlimPajama~\citep{Soboleva:23}. Our evaluation is multi-faceted: we first analyze each model’s ability to extrapolate to longer contextual sequences; next, we measure validation perplexity and assess performance in downstream tasks~\citep{Sutawika:24}; then, we test each method across 571 textual domains using the PALOMA language benchmark suite~\citep{Magnusson:23arxivshort}; finally, we revisit scaling characteristics, now with datasets 20 times larger.

\label{sec:xlstm_blocks_notation}
Throughout all experimental setups, we denote xLSTM[$a$:$b$] as representing the proportion $a/b$ between mLSTM-based and sLSTM-based xLSTM blocks. As an illustration, xLSTM[7:1] indicates that among eight blocks, seven consist of mLSTM-based blocks while one is sLSTM-based. If the total number of blocks is 48, this configuration corresponds to 42 mLSTM-based and 6 sLSTM-based blocks. Additionally, for every experiment, we implement pre up-projection blocks for mLSTM and post up-projection blocks for sLSTM.

\subsection{Synthetic Tasks and Long Range Arena}
\label{sec:ExpSyntethic}
Initially, we examine how well the xLSTM's exponential gating combined with memory mixing operates on formal language benchmarks~\citep{Deletang:23}. Next, we evaluate xLSTM’s recently introduced matrix memory approach using the Multi-Query Associative Recall task~\citep{Arora:23arxiv}. Lastly, we analyze xLSTM's capability for handling extensive sequences by testing it on the Long Range Arena~\citep{Tay:21}.

\paragraph{Evaluation of xLSTM's Exponential Gating Combined with Memory Mixing.} We investigate the performance of xLSTM augmented with exponential gating and memory mixing, which are hypothesized to facilitate effective state tracking~\citep{Merrill:24, Merrill:23}. Building on the formal language benchmarks from \citet{Deletang:23}, we adapt and expand them to include multi-length sequences in the training phase, supporting the evaluation of length extrapolation capabilities. For thorough task details and further empirical results, see Appendix~\ref{sec:appExpSynthetic-formalLang}. We benchmark xLSTM against a diverse range of architectures, encompassing Transformers, State Space Models, and conventional Recurrent Neural Networks. Task-relevant token accuracies are assessed, with all results normalized between 0 (chance) and 1 (optimal). We carry out comparisons involving 2-block model variants, including xLSTM[0:1] (representing sLSTM only), xLSTM[1:0] (mLSTM only), xLSTM[1:1], as well as Llama, Mamba, RWKV, Retention, Hyena, LSTM, and LSTM integrated within Transformer blocks (denoted as LSTM (Block)). Figure~\ref{fig:formal-main} presents the outcomes of these comparisons. Models lacking memory mixing—such as basic Transformers and State Space Models—are observed to be incapable of addressing certain state-dependent tasks, like those defined by regular grammars in the parity task. This aligns with prior observations indicating the inherent limitations of Transformers and State Space Models relative to RNNs in such contexts~\citep{Merrill:24, Merrill:23, Deletang:23}.

\paragraph{Test of xLSTM's Memory Capacities on Associative Recall Tasks.} In this study, we evaluate the newly introduced matrix memory component of xLSTM by assessing its ability to store information on the Multi-Query Associative Recall benchmark~\citep{Arora:23arxiv}. In this task, each sequence is composed of randomly selected key-value pairs from a comprehensive vocabulary; these associations must be retained for subsequent retrieval. To increase the challenge relative to the original setup, we expand the number of key-value pairs up to 256 and push the context length to a maximum of 2048. This setup thus provides a more rigorous assessment of the memory capabilities of the competing models. Our comparative analysis includes 2-block configurations of Llama, Mamba, RWKV-5, RWKV-6, as well as xLSTM[1:1] and xLSTM[1:0]. We determine each model's effectiveness based on its accuracy in recalling the stored associations. Since Transformers like Llama possess memory capacity that grows exponentially with respect to the coding dimension~\citep{Ramsauer:21}, they serve as the reference point for this evaluation. Figure~\ref{fig:mqar-main} displays the main findings. Notably, xLSTM[1:1] achieves superior memory performance compared to all non-Transformer baselines, including in compact model variants. Contrary to intuition, integrating the sLSTM block actually enhances, rather than restricts, memory, as shown most clearly on tasks involving 256 key-value pairs. Further experimental outcomes are included in Appendix~\ref{sec:appExpSynthetic-mqar}, where extrapolation analyses show that the extended memory capacities of xLSTM permit it to generalize to sequence lengths beyond those encountered during training.

\paragraph{Test of xLSTM's Long Context Capabilities on Long Range Arena.} To evaluate how effectively xLSTM processes extended sequences and substantial contexts, we benchmark various approaches on Long Range Arena~\citep{Tay:21}. Across every task, xLSTM reliably achieves high performance, indicating that its architecture excels at managing the complexities associated with long context scenarios.

For more details, see Appendix~\ref{sec:appExpSynthetic-lra}.

\subsection{Method Comparison and Ablation Study}
\label{sec:ExpComparison}

In order to investigate the central inquiry of this study—specifically, how the newly proposed LSTM variants perform when applied at scale in language modeling tasks—we conduct training of xLSTMs, Transformers, State Space Models, among additional techniques, utilizing 15B tokens sourced from SlimPajama within a consistent auto-regressive framework. We then evaluate all models on the validation dataset and carry out ablation experiments for the xLSTMs.

\paragraph{Comparative Analysis of xLSTM and Baseline Approaches.} To provide a comprehensive evaluation, we train each model on a dataset comprising 15B tokens from SlimPajama~\citep{Soboleva:23}. Perplexity on the validation split serves as the principal metric for assessment. The study encompasses a variety of approaches: xLSTM (the proposed architecture), GPT-3 (Transformer-based)~\citep{Brown:20short}, Llama (Transformer-based)~\citep{Touvron:23arxiv}, H3 (SSM)~\citep{Fu:23}, Mamba (SSM)~\citep{Gu:24arxiv}, RWKV-4, RWKV-5, RWKV-6 (all RNNs)~\citep{Peng:23arxivshort, Peng:24arxivshort}, GLA (a linear Transformer)~\citep{Yang:23arxiv}, HGRN and HGRN2 (RNNs)~\citep{Qin:23, Qin:24arxiv}, RetNet (linear Transformer)~\citep{Sun:23arxiv}, Hyena (linear Transformer)~\citep{Poli:23}, as well as the xLSTM[1:0] and xLSTM[7:1] variants (see Section~\ref{sec:xlstm_blocks_notation}). 
Models were trained in mixed-precision mode; however, RWKV-5, RWKV-6, GLA, and HGRN2 did not utilize PyTorch’s automatic mixed-precision feature (refer also to Appendix Section~\ref{sec:appGenTrainProc}).
The examined methodologies fall into three broad categories: (a) Transformers, (b) State Space Models (SSMs), and (c) Recurrent Neural Networks (RNNs) with linear Transformers included. The latter adopt a linear formulation in lieu of the standard attention mechanism found in conventional Transformers. All models are configured to be parameter-equivalent to a 350M-parameter GPT-3, which is achieved by setting the embedding dimension to 1024 with 24 residual blocks. Notably, GPT-3 differs by tying weights between the token and output embeddings, leading to a modestly reduced total parameter count. According to the results in Table~\ref{tab:spaj15b_model_results}, xLSTM delivers superior validation perplexity compared to the existing models. Additional experimental specifics are detailed in Appendix~\ref{sec:appExpComparison}. 
Figure~\ref{fig:temp_sclaw15B} illustrates the scaling trends for this set of experiments, suggesting that xLSTM maintains its competitive advantage as model size increases.

\paragraph{Ablation Studies.}
As shown in Table~\ref{tab:spaj15b_model_results} and Figure~\ref{fig:temp_sclaw15B}, xLSTM delivers strong language modeling performance when trained on 15B tokens from SlimPajama. To systematically evaluate the contributions of individual modifications, we incrementally transform a standard LSTM into an xLSTM by sequentially incorporating new components. The process begins by embedding LSTM layers within pre-LayerNorm residual architectures. Next, we expand this to include a post up-projection module. In the final step, we introduce exponential gating along with matrix memory mechanisms.

The results are shown in Table~\ref{tab:ablstudies} (top). 

Ablation experiments indicate that the significant gains observed are primarily due to the exponential gating and the inclusion of matrix memory. Recognizing the critical role of gating in both RNNs and State Space Models, we further explore various gating strategies by systematically ablating them. The results presented in Table~\ref{tab:ablstudies} (bottom) demonstrate that allowing each gate to be trainable and responsive to the input yields progressively better outcomes. Further analyses concerning the backbone elements are provided in Appendix~\ref{sec:appAblDetails}.

\subsection{xLSTM as Large Language Model}
\label{sec:ExpLanguage}

To conclude, we conduct extensive language modeling experiments at scale, evaluating the effectiveness of xLSTM in the context of large language models. The training data size is increased, with models being trained on 300B tokens sourced from SlimPajama. Notably, this matches the token counts used in models such as Mamba~\citep{Gu:24arxiv} and Griffin~\citep{De:24arxiv}.

Our evaluation benchmarks xLSTM against RWKV-4, Llama, and Mamba, each serving as an exemplar from the method categories outlined in Section~\ref{sec:ExpComparison}. RWKV-4 is chosen to represent RNN-based approaches, as satisfactory training precision configurations for RWKV-5, RWKV-6, and HGRN2 were only discovered subsequent to the commencement of the 300B token training experiments (refer to Appendix~\ref{sec:appGenTrainProc} for details).

We implement and train a spectrum of model sizes (125M, 350M, 760M, 1.3B) and evaluate them for both their ability to extrapolate to longer sequences and their accuracy on validation datasets. Further, the models are examined on a variety of downstream tasks, tested across 571 domains of the PALOMA benchmark to analyze language modeling performance, and assessed with respect to their scaling law dynamics.

\paragraph{Sequence Length Extrapolation.}
\label{sec:spaj300B_ctx_extrapolation}
To begin, we investigate the ability of 1.3B-parameter versions of xLSTM, RWKV-4, Llama, and Mamba to generalize to longer sequences. Each model is trained using a context length of 2048 tokens and later evaluated on extended contexts, scaling up to 16384 tokens. Results are depicted in Figure~\ref{fig:spaj300B_seq_extrapolation}. Unlike the other approaches, xLSTM consistently exhibits low perplexity over increased sequence lengths.

\paragraph{Validation Perplexity and Downstream Tasks.}
\label{sec:spaj_downstream_eval}
Next, across varying model sizes, we assess xLSTM, RWKV-4, Llama, and Mamba on the SlimPajama validation set for next token prediction, as well as on a suite of downstream tasks targeting common sense reasoning skills.

The third column in 
Table~\ref{tab:lmeval} presents the validation set perplexities achieved by various approaches. For every model size considered, xLSTM[1:0] and xLSTM[7:1] yield the lowest perplexity on the validation set, indicating their superior performance. Performance on downstream tasks is detailed in the remaining columns of 
Table~\ref{tab:lmeval}. xLSTM outperforms all other methods on most tasks and across all model sizes, except in select cases for the ARC task, where Mamba performs better. Additional discussion and specifics can be found in Appendix~\ref{sec:appExpLanguage}.

\paragraph{Performance on PALOMA Language Tasks.} As a third point, we evaluated xLSTM, RWKV-4, Llama, and Mamba models across all size variants by assessing their next token prediction ability on the PALOMA language tasks~\citep{Magnusson:23arxivshort}. This evaluation utilized perplexity scores on next token predictions within 571 distinct textual domains, covering a spectrum that includes sources such as nytimes.com and Reddit's r/depression.

Table~\ref{tab:ppleval} presents the results, with perplexity values organized into two categories: general language modeling tasks (shown in the first seven columns) and specific domain benchmarks (displayed in the final five columns). For these language tasks, xLSTM[1:0] consistently yields lower perplexity than its counterpart xLSTM[7:1].

The xLSTM[1:0] model achieves lower perplexity than Mamba in 568 out of 571 text domains (99.5\%), outperforms Llama in 486 domains (85.1\%), and surpasses RWKV-4 in 570 domains (99.8\%). Additional data are available in Appendix~\ref{sec:appExpLanguage}.

\paragraph{Scaling Laws.} Next, we investigate the power-law scaling patterns, which enable predictions of model performance as their size increases~\citep{Kaplan:20,Brown:20short}. 
The scaling results are shown in Figure~\ref{fig:temp_sclaw300B}. Although each model exhibits comparable scaling trends, they differ in their baseline performance levels. Among the evaluated architectures, RWKV-4 demonstrates the lowest performance, with Llama and Mamba outperforming it. The xLSTM model surpasses Mamba, maintaining roughly the same advantage that Mamba holds over Llama. These scaling curves suggest that as models grow, xLSTM maintains a competitive edge relative to both Transformer-based and State-Space models.

\textbf{Generation Times and Maximal Throughput.} We present measurements of text generation latency in Figure~\ref{fig:inference_time_generation} and report peak throughput results in Figure~\ref{fig:inference_time_decoding_throughput} (left) for the xLSTM models at the 1.3B parameter count. For comparison, we evaluate similarly sized implementations of Mamba, Llama, and RWKV from HuggingFace; the Llama baselines make use of a static key-value cache. Notably, during the time of these experiments, neither full cache compilation for Transformer models nor the usage of \texttt{torch.compile} on the Mamba model was functional. In generation benchmarks, all models are assessed with batch size set to 1 and pre-fill set to 16, ensuring conditions that advantage Transformer architectures.

As shown in Figure~\ref{fig:inference_time_generation}, both xLSTM and other recurrent methods (Mamba and RWKV-4) display linear time growth, as opposed to the quadratic complexity observed in Llama. To assess decoding throughput, we further vary batch size and prefill values for the Llama model. Results in Figure~\ref{fig:inference_time_decoding_throughput} (right) indicate that xLSTM supports substantially larger batch sizes than Llama, thanks to its fixed memory footprint; as a consequence, xLSTM achieves superior throughput.

\section{Limitations}

(i) Unlike mLSTM, the sLSTM's approach to memory mixing prevents parallel computation, which in turn inhibits efficient parallelized implementations. Despite this, we engineered a specialized CUDA kernel for sLSTM that performs at less than twice the runtime of our parallel mLSTM kernel. (ii) It should be noted that the existing mLSTM CUDA kernels lack optimization, resulting in performance roughly four times slower than FlashAttention or Mamba’s scan mechanism. Performance gains are possible through adoption of strategies similar to those in FlashAttention. (iii) mLSTM's matrix memory necessitates significant computational effort, as operations must be conducted on $d \times d$ matrices. Nevertheless, because both memory update and retrieval are parameter-free and can exploit standard parallel matrix operations, the additional runtime incurred from this complexity is relatively small. (iv) Proper selection of forget gate initialization is essential. (v) As the matrix memory size in mLSTM does not depend on sequence length, extending the sequence may place excessive demand on memory with long contexts. However, this issue is not evident for contexts up to 16k tokens, see Section~\ref{sec:spaj300B_ctx_extrapolation}. (vi) Owing to substantial computation requirements in large language model experiments, we have not undertaken comprehensive architectural or hyperparameter tuning, especially regarding larger xLSTM variants. We believe that xLSTM's true capabilities would be realized through rigorous optimization procedures.

\section{Conclusion}
We have addressed our core inquiry: What advancements in language modeling can be achieved by scaling LSTM architectures to billions of parameters? Our findings show that this approach reaches performance levels on par with established methodologies such as Transformers and State Space Models. By introducing exponential gating with memory mixing and a novel memory structure, we extend LSTM into xLSTM. The xLSTM variants demonstrate strong performance in language modeling tasks, exhibiting competitiveness with current leading approaches including Transformers and State Space Models. Analysis of scaling laws suggests that expanding xLSTM models may position them as viable alternatives to large Transformer-based language models. Furthermore, xLSTM opens promising opportunities for application in domains such as Reinforcement Learning, Time Series Prediction, and physical system modeling.

\end{document}